\documentclass[12pt]{article}
\usepackage[utf8]{inputenc}
\usepackage[T1]{fontenc}
\usepackage{hyperref}
\title{Chumash Language Resources for K-12 Curriculum Development}
\author{Compiled by Claude Sonnet 5.5 (claude-sonnet-5-5) for Andr\'as Kornai}
\date{}
\begin{document}
\maketitle

\section{Introduction}
Chumash (ISO 639-3 chs) is a group of Indigenous languages of southern California. The last known native speakers died in the 20th century (the last Barbareño speaker in 1965), and current work is revival based on archival notes, chiefly those of J.\,P.\ Harrington. The searches mostly concern Barbareño/Ventureño and Ineseño (Samala); I did not verify which varieties ISO chs covers.

\section{Primary Educational Resources}
\begin{itemize}
\item Samala is reported taught at two local elementary schools, and the Santa Ynez Band runs an immersion-based apprentice program. \url{https://www.nps.gov/chis/learn/news/bringing-back-the-samala-chumash-language.htm}
\item Wishtoyo Chumash Village (Malibu) opened its language school in 2010, per a summary of the revival. \url{https://omniglot.com/writing/barbareno.htm}
\end{itemize}

\section{Linguistic and Reference Materials}
\begin{itemize}
\item Samala--English dictionary (Santa Ynez Band with Richard Applegate, 2008; 608 pages, with pronunciation disc). \url{https://www.independent.com/2008/04/24/chumash-dictionary-breathes-life-into-moribund-language/}
\item Online Barbareño/Ventureño Chumash dictionary hosted by CSU Channel Islands (\emph{mitsqanaqan}, by Timothy Henry-Rodriguez). \url{https://www.csuci.edu/news/channel/spring-2024/chumashlanguage.htm}
\item Local coverage of the online dictionary. \url{https://www.conejovalleyguide.com/local-buzz/csu-channel-islands-hosts-a-new-online-chumash-language-dictionary}
\end{itemize}

\section{Religious and Community Resources}
Santa Ynez Band of Chumash Indians culture materials. \url{https://chumash.gov/s/Culture.pdf}

\section{Digital and Audio Resources}
The online dictionary above is the main digital resource found; the 2008 dictionary includes a pronunciation disc.

\section{Curriculum Development Context}
Revival rests on archival records and newly trained learners; no complete K-12 curriculum was found in the searches. Any curriculum should be developed with the tribal language programs named above.

\end{document}